% !TeX root = ../../main.tex
\section{\bilingual{Formalization and Branching Properties}{形式化与分支性质}}
\label{section:formalization-and-branching-properties}

\bilingual{
    Consider an implicit probability space $(\Omega,\mathcal F,\mathbb P)$, on which all random objects in the following are defined. For the value space of vertices of trees, we adopt the terminology of Ulam--Harris--Neveu \mycite{Neveu1986}, which is defined as follows
    \begin{equation*}
        \mathbb V:=\bigcup_{n=0}^\infty\mathbb N^n,
    \end{equation*}
    where $\mathbb N^0:=\{\emptyset\}$, $\emptyset$ is the root of the tree, $\mathbb N^n$ is the $n$-dimensional natural number space. The length function $|\cdot|:\mathbb V\to\mathbb N$ defined on it is called the generation.  Denote $u_i$ as the ancestor of $u\in\mathbb V$ at generation $i$. A {\itshape tree} (tree) or {\itshape rooted labeled tree} is a subset $T$ of $\mathbb V$ that satisfies the following conditions:
    \begin{itemize}
        \item $\emptyset\in T$;
        \item If $u\in T$ and $u=vw$ where $v,w\in\mathbb V$, $vw$ denotes the concatenation of the two, then $v\in T$;
        \item If $uj\in T$, $j\in\mathbb N$, then $uj'\in T$ for any $j'\in[\![1,j]\!]$.
    \end{itemize}
    The set of all trees is denoted by $\mathbb T$. A {\itshape finite tree} is a tree with a finite number of vertices. A {\itshape locally finite tree} is a tree whose truncation at each generation is finite, and the set of all locally finite trees is denoted by $\mathbb T_0$. The space $\mathbb T$ naturally defines a metric
    \begin{equation*}
        d(T_1,T_2):=e^{-\inf\{n\ge 0:T_1\cap\mathbb N^n\neq T_2\cap\mathbb N^n\}},
    \end{equation*}
    where $T_1,T_2\in\mathbb T$, but it is only complete on the space of locally finite trees $\mathbb T_0$, whose countable dense set is the set of all finite trees. For the $\sigma$-algebra on $\mathbb T$, we use the coarser $\sigma$-algebra $\mathcal G_{\mathbb T}$ generated by the cylinder sets
    \begin{equation*}
        \left\{
            T\in\mathbb T: u\in T
        \right\},\quad u\in\mathbb V
    \end{equation*}
    Denote the random tree taking values in $(\mathbb T,\mathcal G_{\mathbb T})$ as $\mathcal T$.
}{
    考虑一个隐含的概率空间$(\Omega,\mathcal F,\mathbb P)$, 后文的所有随机对象都定义在这个概率空间上. 对于树上节点的取值空间, 采用 Ulam--Harris--Neveu \mycite{Neveu1986} 的术语, 其定义如下
    \begin{equation*}
        \mathbb V:=\bigcup_{n=0}^\infty\mathbb N^n,
    \end{equation*}
    其中 $\mathbb N^0:=\{\emptyset\}$, $\emptyset$ 是树的根节点, $\mathbb N^n$ 是 $n$ 维的自然数空间. 其上的长度函数$|\cdot|:\mathbb V\to\mathbb N$称为代数. 记$u_i$为$u\in\mathbb V$的第$i$代的祖先. 所谓{\kaishu 树} (tree) 或{\kaishu 有根标签树} (rooted labeled tree), 是指 $\mathbb V$ 的一个子集 $T$ 满足如下条件:
    \begin{itemize}
        \item $\emptyset\in T$;
        \item 若 $u\in T$ 且 $u=vw$ 其中 $v,w\in\mathbb V$, $vw$指两者的连接, 则 $v\in T$;
        \item 若 $uj\in T$, $j\in\mathbb N$, 则 $uj'\in T$ 对任意 $j'\in[\![1,j]\!]$ 成立.
    \end{itemize}
    树的全体记作$\mathbb T$. 所谓{\kaishu 有限树} (finite tree), 是指节点数有限的树. 所谓{\kaishu 局部有限树} (locally finite tree), 是指每一层的截断都有限的树, 全体记作$\mathbb T_0$. 一般采用$\mathbb T_0$空间作为考察对象, 但本文需要考虑无限分支的分支律, 所得到的树可能不满足局部有限的条件, 因此也需要考虑$\mathbb T$空间. 树空间$\mathbb T$上自然定义有一个度量
    \begin{equation*}
        d(T_1,T_2):=e^{-\inf\{n\ge 0:T_1\cap\mathbb N^n\neq T_2\cap\mathbb N^n\}},
    \end{equation*}
    其中 $T_1,T_2\in\mathbb T$, 但它仅在局部有限树空间$\mathbb T_0$上完备, 其可数稠密集为有限树全体. 对于$\mathbb T$上的$\sigma$-代数, 转而采用柱集
    \begin{equation*}
        \left\{
            T\in\mathbb T: u\in T
        \right\},\quad u\in\mathbb V
    \end{equation*}
    生成的更粗的$\sigma$-代数$\mathcal G_{\mathbb T}$. 记取值于$(\mathbb T,\mathcal G_{\mathbb T})$的随机树为$\mathcal T$.
}

\bilingual{
    Sometimes, we also consider a family of trees indexed by $\mathbb I$ (with a common root $\emptyset$), whose whole is $\mathbb T^{\mathbb I}$, and its $\sigma$-algebra is $\mathcal G_{\mathbb T}^{\otimes\mathbb I}$, which is also generated by cylinder sets. Denote the random tree family taking values in $(\mathbb T^{\mathbb I},\mathcal G_{\mathbb T}^{\otimes\mathbb I})$ as $\boldsymbol{\mathcal T}$.  Moreover, for some $(u,i)\in\boldsymbol T$, denote $\boldsymbol T_{u,i}$ as the subtree rooted at $u$, which is defined as
    \begin{equation*}
        \boldsymbol T_{u,i}:=\{v\in\mathbb V:uv\in\boldsymbol T(i)\}.
    \end{equation*}
}{
    有时也会考虑由$\mathbb I$索引的一族 (有共同根$\emptyset$的) 树$\boldsymbol T$, 其全体为$\mathbb T^{\mathbb I}$, 其$\sigma$-代数为$\mathcal G_{\mathbb T}^{\otimes\mathbb I}$, 事实上也是由柱集生成的. 记取值于$(\mathbb T^{\mathbb I},\mathcal G_{\mathbb T}^{\otimes\mathbb I})$的随机树族为$\boldsymbol{\mathcal T}$. 另外, 对于某个$(u,i)\in\boldsymbol T$, 记$\boldsymbol T_{u,i}$为以$u$为根节点的子树, 其定义为
    \begin{equation*}
        \boldsymbol T_{u,i}:=\{v\in\mathbb V:uv\in\boldsymbol T(i)\}.
    \end{equation*}
}

\bilingual{
    The so-called {\itshape marked tree} is defined as follows: it is a tree $T\in\mathbb T$ equipped with a marking function $V:T\to\mathbb R$; denoted as $(T,V)$ or simply $V$. We require that $V(\emptyset)=0$. The set of all marked trees is denoted by
    \begin{equation*}
        \mathbb M:=\{(T,V):T\in\mathbb T,\,V:T\to\mathbb R\}.
    \end{equation*}
    The $\sigma$-algebra $\mathcal G_{\mathbb M}$ on it is generated by the cylinder sets
    \begin{equation*}
        \left\{(T,V)\in\mathbb M: u\in T,\,V(u)\in A\right\},\quad u\in\mathbb V,\,A\in\mathcal B(\mathbb R).
    \end{equation*}
    Denote the random marked tree taking values in $(\mathbb M,\mathcal G_{\mathbb M})$ as $(\mathcal T,\mathcal V)$, sometimes simply denoted as $\mathcal V$, and called a {\itshape branching random walk} (BRW).
}{
    所谓{\kaishu 标记树} (marked tree), 是指在某棵树$T\in\mathbb T$的基础上, 还定义了一个标记函数 $V:T\to\mathbb R$; 记作$(T,V)$或直接记为$V$. 要求$V(\varnothing)=0$. 标记树全体记为
    \begin{equation*}
        \mathbb M:=\{(T,V):T\in\mathbb T,\,V:T\to\mathbb R\}.
    \end{equation*}
    其上的$\sigma$-代数$\mathcal G_{\mathbb M}$由柱集
    \begin{equation*}
        \left\{(T,V)\in\mathbb M: u\in T,\,V(u)\in A\right\},\quad u\in\mathbb V,\,A\in\mathcal B(\mathbb R)
    \end{equation*}
    生成, 记取值于$(\mathbb M,\mathcal G_{\mathbb M})$的随机标记树为$(\mathcal T,\mathcal V)$, 有时也直接记为$\mathcal V$, 称为{\kaishu 分支随机游走} (branching random walk, BRW).
}

\bilingual{
    Similarly, we also consider a family of marked trees indexed by $\mathbb I$. For a family of trees $\boldsymbol T\in\mathbb T^{\mathbb I}$, the marking function is still denoted as $V$, but its domain is considered as
    $
        \left\{
            (u,i):u\in\boldsymbol T(i),\,i\in\mathbb I
        \right\}
    $. We require that $V(\varnothing,i)=0$. Such marked trees are denoted as $(\boldsymbol T,V)$ or simply $V$, and the set of all such marked trees is denoted by
    \begin{equation*}
        \mathbb M^{\mathbb I}:=\{(\boldsymbol T,V):\boldsymbol T\in\mathbb T^{\mathbb I},\,V:\{(u,i):u\in\boldsymbol T(i),\,i\in\mathbb I\}\to\mathbb R\}.
    \end{equation*}
    Its $\sigma$-algebra $\mathcal G_{\mathbb M}^{\otimes\mathbb I}$ is generated by the cylinder sets
    \begin{equation*}
        \left\{(\boldsymbol T,V)\in\mathbb M^{\mathbb I}: u\in\boldsymbol T(i),\,V(u,i)\in A\right\},\quad u\in\mathbb V,\,i\in\mathbb I,\,A\in\mathcal B(\mathbb R)
    \end{equation*}
    Denote the random marked tree family taking values in $(\mathbb M^{\mathbb I},\mathcal G_{\mathbb M}^{\otimes\mathbb I})$ as $(\boldsymbol{\mathcal T},\mathcal V)$, sometimes simply denoted as $\mathcal V$.
}{
    同样地, 也会考虑由$\mathbb I$索引的一族标记树. 对于一族树$\boldsymbol T\in\mathbb T^{\mathbb I}$, 标记函数仍记作$V$, 但其定义域视作
    $
        \left\{
            (u,i):u\in\boldsymbol T(i),\,i\in\mathbb I
        \right\}
    $.
    要求$V(\varnothing,i)=0$. 这样的标记树族记为$(\boldsymbol T,V)$或直接记为$V$, 其全体记为
    \begin{equation*}
        \mathbb M^{\mathbb I}:=\{(\boldsymbol T,V):\boldsymbol T\in\mathbb T^{\mathbb I},\,V:\{(u,i):u\in\boldsymbol T(i),\,i\in\mathbb I\}\to\mathbb R\}.
    \end{equation*}
    其$\sigma$-代数$\mathcal G_{\mathbb M}^{\otimes\mathbb I}$由柱集
    \begin{equation*}
        \left\{(\boldsymbol T,V)\in\mathbb M^{\mathbb I}: u\in\boldsymbol T(i),\,V(u,i)\in A\right\},\quad u\in\mathbb V,\,i\in\mathbb I,\,A\in\mathcal B(\mathbb R)
    \end{equation*}
    生成, 记取值于$(\mathbb M^{\mathbb I},\mathcal G_{\mathbb M}^{\otimes\mathbb I})$的随机标记树族为$(\boldsymbol{\mathcal T},\mathcal V)$, 或直接记为$\mathcal V$.
}

\bilingual{
    Now consider the countable-length {\itshape real vector space} $\mathbb R_{\mathrm V}$, which is defined as
    \begin{equation*}
        \mathbb R_{\mathrm V}:=\mathbb R^{\mathbb N}\cup\bigcup_{n=0}^\infty\mathbb R^n,
    \end{equation*}
    where $\mathbb R^0:=\{\emptyset\}$, $\emptyset$ is the zero vector of zero dimension. The length function $\#:\mathbb R_{\mathrm V}\to\mathbb N\cup\{\infty\}$ and the component function $(\cdot)^i:\mathbb R_{\mathrm V}\to\mathbb R$ are defined on it. The $\sigma$-algebra $\mathcal G_{\mathbb R_{\mathrm V}}$ on it is generated by the family of sets
    \begin{equation*}
        \left\{
            x\in\mathbb R_{\mathrm V}: \#x\ge i,\,x^i\in A
        \right\},\quad i\in\mathbb N,\,A\in\mathcal B(\mathbb R)
    \end{equation*}
    The so-called branching law $\Xi$ is a random vector taking values in $(\mathbb R_{\mathrm V},\mathcal G_{\mathbb R_{\mathrm V}})$, whose component $\Xi^i$ represents the displacement of the $i$-th offspring, and its length $\#\Xi$ represents the number of offspring. In the following, under certain assumptions, if there exists a leftmost component, we will denote $\Xi_i$ as the $i$-th leftmost (smallest) component of $\Xi$.
}{
    现在考虑包含可数长度向量的{\kaishu 实向量空间} $\mathbb R_{\mathrm V}$, 即
    \begin{equation*}
        \mathbb R_{\mathrm V}:=\mathbb R^{\mathbb N}\cup\bigcup_{n=0}^\infty\mathbb R^n,
    \end{equation*}
    其中 $\mathbb R^0:=\{\emptyset\}$, $\emptyset$ 是零维的零向量. 其上定义有长度函数$\#:\mathbb R_{\mathrm V}\to\mathbb N\cup\{\infty\}$ 与分量函数$(\cdot)^i:\mathbb R_{\mathrm V}\to\mathbb R$. 其上的$\sigma$-代数$\mathcal G_{\mathbb R_{\mathrm V}}$由如下集族
    \begin{equation*}
        \left\{
            x\in\mathbb R_{\mathrm V}: \#x\ge i,\,x^i\in A
        \right\},\quad i\in\mathbb N,\,A\in\mathcal B(\mathbb R)
    \end{equation*}
    生成. 所谓分支律$\Xi$, 就是一个取值于$(\mathbb R_{\mathrm V},\mathcal G_{\mathbb R_{\mathrm V}})$的随机向量, 其分量$\Xi^i$表示第$i$个子代的位移, 其长度$\#\Xi$表示子代的数量. 下文中, 在一定的假设下, 若存在最左侧分量, 将用$\Xi_i$表示$\Xi$的最左侧 (最小的) 第$i$个分量.
}

\bilingual{
    Given a branching law $\Xi$, consider an index set $\mathbb I$, then a family of i.i.d. branching laws $\{\Xi_{u,i}\}_{u\in\mathbb V,\,i\in\mathbb I}$ determines a random marked tree family $(\boldsymbol{\mathcal T},\mathcal V)$, where $\boldsymbol{\mathcal T}(i)$ is the Galton--Watson tree generated by this family of branching laws in sequence, $\mathcal V(u,i)$ is the displacement of vertex $u$ on tree $\boldsymbol{\mathcal T}(i)$, satisfying the following relation
    \begin{equation*}
        \Xi_{u,i}=\mathcal V(u,i)-\mathcal V(u_{|u|-1},i),
        \quad i\in\mathbb I,\,u\in\boldsymbol{\mathcal T}(i)-\{\emptyset\}.
    \end{equation*}
    In the absence of ambiguity, the index $i$ is usually omitted. Moreover, denote the set of all particles in the $n$-th generation as
    \begin{equation*}
        \mathcal Z_n:=\{(u,i):u\in\boldsymbol{\mathcal T}(i),\,i\in\mathbb I,\,|u|=n\},
    \end{equation*}
    and its quantity as $\#\mathcal Z_n$. Also denote the infimum of the particle positions in the $n$-th generation as
    \begin{equation*}
        \mathcal M_n:=\inf\{\mathcal V(u,i):(u,i)\in\mathcal Z_n\}.
    \end{equation*}
    In the following, under certain assumptions, this infimum exists and is a minimum.
}{
    给定分支律 $\Xi$ 与索引集 $\mathbb I$. 对每个 $(u,i)\in\mathbb V\times\mathbb I$ 取一个独立同分布的分支律 $\Xi_{u,i}$, 由此逐代生成随机标记树族 $(\boldsymbol{\mathcal T},\mathcal V)$. 每棵树 $\boldsymbol{\mathcal T}(i)$ 都是 Galton--Watson 树, 节点位移满足
    \begin{equation*}
        \Xi_{u,i}=\mathcal V(u,i)-\mathcal V(u_{|u|-1},i),
        \quad i\in\mathbb I,\,u\in\boldsymbol{\mathcal T}(i)-\{\emptyset\}.
    \end{equation*}
    在不造成歧义的情况下, 通常省略索引$i$. 另外记第$n$代的粒子全体为
    \begin{equation*}
        \mathcal Z_n:=\{(u,i):u\in\boldsymbol{\mathcal T}(i),\,i\in\mathbb I,\,|u|=n\},
    \end{equation*}
    其数量记为$\#\mathcal Z_n$. 还记第$n$代的粒子位置的下确界为
    \begin{equation*}
        \mathcal M_n:=\inf\{\mathcal V(u,i):(u,i)\in\mathcal Z_n\}.
    \end{equation*}
    下文中, 在一定的假设下, 此下确界存在且为最小值.
}

\bilingual{
    Now consider the generation filtration generated by the family of branching laws $(\Xi_{u,i})$ on the implicit probability space. Define
    \begin{equation*}
        \mathcal F_n:=\begin{cases}
            \sigma\left(
                \Xi_{u,i}:u\in\mathbb V,\,i\in\mathbb I,\,|u|\le n-1
            \right),& n\in\mathbb Z_+,\\
            \{\emptyset,\Omega\},& n=0,
        \end{cases}
    \end{equation*}
    which is also equal to
    $
        \sigma\left(
            \mathcal V(u,i):u\in\mathbb V,\,i\in\mathbb I,\,|u|\le n
        \right)
    $.
}{
    现在考虑分支律族$(\Xi_{u,i})$在隐含概率空间上生成的逐代域流. 定义
    \begin{equation*}
        \mathcal F_n:=\begin{cases}
            \sigma\left(
                \Xi_{u,i}:u\in\mathbb V,\,i\in\mathbb I,\,|u|\le n-1
            \right),& n\in\mathbb Z_+,\\
            \{\emptyset,\Omega\},& n=0,
        \end{cases}
    \end{equation*}
    它也等于
    $
        \sigma\left(
            \mathcal V(u,i):u\in\mathbb V,\,i\in\mathbb I,\,|u|\le n
        \right)
    $.
}

\bilingual{
    The following branching properties are important tools in this thesis.
}{
    下面的分支性质是本文的重要工具.
}

\begin{lemma}[\bilingual{Branching Property}{分支性质}]\it\label{lemma:branching-property}
    \bilingual{
        Given a branching law $\Xi$. With the same notation as above. For any $n\in\mathbb Z_+$, conditional on $\mathcal F_n$ and $\#\mathcal Z_n\ge 1$, the family of sub-marked trees rooted at $\{(u,i):u\in\mathcal Z_n\}$, i.e., $\{(\boldsymbol{\mathcal T}\!\!{}_{u,i},\,v\mapsto\mathcal V(uv,i)-\mathcal V(u,i)):(u,i)\in\mathcal Z_n\}$, is a family of i.i.d. random marked trees, whose distribution is the same as some $(\boldsymbol{\mathcal T}(i_0),\mathcal V(\cdot,i_0))$, $i_0\in\mathbb I$.
    }{
        给定分支律$\Xi$, 以及由它决定的随机标记树族$(\boldsymbol{\mathcal T},\mathcal V)$, 则对于任意 $n\in\mathbb Z_+$, 条件于 $\mathcal F_n$ 以及 $\#\mathcal Z_n\ge 1$, 以 $\{(u,i):u\in\mathcal Z_n\}$ 为根节点的子标记树族$\{(\boldsymbol{\mathcal T}\!\!{}_{u,i},\,v\mapsto\mathcal V(uv,i)-\mathcal V(u,i)):(u,i)\in\mathcal Z_n\}$ 是一族独立同分布的随机标记树, 其分布与某个$(\boldsymbol{\mathcal T}(i_0),\mathcal V(\cdot,i_0))$, $i_0\in\mathbb I$相同.
    }
\end{lemma}

\begin{lemma}[\bilingual{Branching Property at Stopping Times}{停时处的分支性质}]\it
    \label{lemma:branching-property-stopping-time}
    \bilingual{
        Given a branching law $\Xi$. With the same notation as above. For an $(\mathcal F_n)$-stopping time $\tau$, conditional on $\mathcal F_\tau$ and $\#\mathcal Z_\tau\ge 1$, the family of sub-marked trees rooted at $\{(u,i):u\in\mathcal Z_\tau\}$, i.e., $\{(\boldsymbol{\mathcal T}\!\!{}_{u,i},\,v\mapsto\mathcal V(uv,i)-\mathcal V(u,i)):(u,i)\in\mathcal Z_\tau\}$, is a family of i.i.d. random marked trees, whose distribution is the same as some $(\boldsymbol{\mathcal T}(i_0),\mathcal V(\cdot,i_0))$, $i_0\in\mathbb I$.
    }{
        给定分支律 $\Xi$. 记号同上. 对于$(\mathcal F_n)$-停时$\tau$, 条件于 $\mathcal F_\tau$ 以及 $\#\mathcal Z_\tau\ge 1$, 以 $\{(u,i):u\in\mathcal Z_\tau\}$ 为根节点的子标记树族$\{(\boldsymbol{\mathcal T}\!\!{}_{u,i},\,v\mapsto\mathcal V(uv,i)-\mathcal V(u,i)):(u,i)\in\mathcal Z_\tau\}$ 是一族独立同分布的随机标记树, 其分布与某个$(\boldsymbol{\mathcal T}(i_0),\mathcal V(\cdot,i_0))$, $i_0\in\mathbb I$相同.
    }
\end{lemma}
\bilingual{
    The above branching properties are special cases of stronger branching properties about {\itshape stopping lines} \mycite{chauvin1986propriete,JAGERS1989183}. A {\itshape line} $\ell$ of a tree is a set of vertices in $\mathbb V$ that are pairwise non-ancestral, the corresponding $\sigma$-algebra $\mathcal G_\ell$ is generated by the branching r.v. at points not in $\ell$ and their descendants, a stopping line is a random line such that any event of being before any deterministic line is measurable.
}{
    \noindent 上述分支性质是更强版本 \mycite{chauvin1986propriete,JAGERS1989183} 的关于{\kaishu 停止线} (stopping lines) 的分支性质的特例. 所谓{\kaishu 线} (line)是指$\mathbb V$中两两无继承关系的点集 $\ell$, 对应的$\sigma$-代数$\mathcal G_\ell$由非$\ell$及其后代处的分支生成, 所谓停止线是指使得其位置在任何确定线前的事件均可测的随机线.
}

\raggedbottom
